\documentclass[11pt,letterpaper]{article}
% KJO_preamble.tex
% Shared LaTeX preamble for all KJO documentation documents.
% \input this file immediately after \documentclass{}.

% --- Encoding and fonts -------------------------------------------------------
\usepackage[utf8]{inputenc}
\usepackage[T1]{fontenc}
\usepackage{lmodern}
\usepackage{microtype}

% --- Page geometry ------------------------------------------------------------
\usepackage[
  letterpaper,
  top=1in, bottom=1in,
  left=1.25in, right=1in
]{geometry}

% --- Core utilities -----------------------------------------------------------
\usepackage{amsmath}
\usepackage{graphicx}
\usepackage{xcolor}
\usepackage{parskip}        % space between paragraphs; no indent
\usepackage{enumitem}
\usepackage{caption}
\usepackage{subcaption}
\usepackage{float}

% --- Tables -------------------------------------------------------------------
\usepackage{booktabs}
\usepackage{longtable}
\usepackage{array}
\usepackage{multirow}

% --- Code listings ------------------------------------------------------------
\usepackage{listings}
\lstset{
  basicstyle=\ttfamily\small,
  keywordstyle=\bfseries\color{KJOblue},
  commentstyle=\itshape\color{gray},
  stringstyle=\color{KJOgreen!70!black},
  breaklines=true,
  frame=single,
  framesep=4pt,
  xleftmargin=6pt,
  xrightmargin=4pt,
  numbers=left,
  numberstyle=\tiny\color{gray},
  numbersep=8pt,
  language=C++,
}

% --- TikZ / PGF ---------------------------------------------------------------
\usepackage{tikz}
\usetikzlibrary{
  arrows.meta,
  shapes.geometric,
  shapes.misc,
  shapes.symbols,
  positioning,
  fit,
  calc,
  backgrounds,
  decorations.pathreplacing,
  decorations.markings,
  matrix,
}
\usepackage{pgfplots}
\pgfplotsset{compat=1.18}

% --- Headers / footers --------------------------------------------------------
\usepackage{fancyhdr}
\usepackage{lastpage}
\pagestyle{fancy}
\fancyhf{}
\renewcommand{\headrulewidth}{0.4pt}
\renewcommand{\footrulewidth}{0.4pt}
% Per-document: \lhead{}, \rhead{}, \cfoot{} are set in the main .tex file.

% --- Section formatting -------------------------------------------------------
\usepackage{titlesec}
\titleformat{\section}{\large\bfseries\color{KJOblue}}{{\thesection}}{0.8em}{}[\titlerule]
\titleformat{\subsection}{\normalsize\bfseries\color{KJOblue!80!black}}{{\thesubsection}}{0.8em}{}
\titleformat{\subsubsection}{\normalsize\bfseries}{{\thesubsubsection}}{0.8em}{}

% --- Hyperlinks ---------------------------------------------------------------
\usepackage[
  colorlinks=true,
  linkcolor=KJOblue,
  urlcolor=KJOblue!70!black,
  citecolor=KJOblue,
  pdfborder={0 0 0},
]{hyperref}

% --- Color palette -----------------------------------------------------------
\definecolor{KJOblue}   {RGB}{ 30,  90, 160}   % RCU
\definecolor{KJOred}    {RGB}{180,  30,  30}   % EMU
\definecolor{KJOgreen}  {RGB}{ 20, 120,  60}   % GSEMU
\definecolor{KJOorange} {RGB}{200, 110,  20}   % GSE / pressurization
\definecolor{KJOgray}   {RGB}{100, 100, 100}   % Buses / links
\definecolor{KJOpurple} {RGB}{110,  40, 140}   % Shared libraries
\definecolor{KJObg}     {RGB}{248, 248, 252}   % Light background for code

% --- Convenience macros -------------------------------------------------------

% Hardware component name (small caps)
\newcommand{\hw}[1]{\textsc{#1}}

% Software module name (monospace)
\newcommand{\sw}[1]{\texttt{#1}}

% Command name (monospace bold)
\newcommand{\cmd}[1]{\texttt{\textbf{#1}}}

% Pin / signal name
\newcommand{\sig}[1]{\texttt{#1}}

% Note / callout box
\usepackage{tcolorbox}
\tcbuselibrary{skins}
\newtcolorbox{notebox}{
  enhanced,
  colback=KJOblue!5,
  colframe=KJOblue!40,
  fonttitle=\bfseries,
  title=Note,
  left=4pt, right=4pt, top=2pt, bottom=2pt,
}
\newtcolorbox{warningbox}{
  enhanced,
  colback=KJOred!5,
  colframe=KJOred!40,
  fonttitle=\bfseries\color{KJOred},
  title=Warning,
  left=4pt, right=4pt, top=2pt, bottom=2pt,
}

% Title page macro: \KJOtitlepage{title}{subtitle}{version}{date}{description}
\newcommand{\KJOtitlepage}[5]{%
  \begin{titlepage}
    \centering
    \vspace*{2cm}
    {\Huge\bfseries\color{KJOblue} #1 \par}
    \vspace{0.5cm}
    {\Large\color{KJOgray} #2 \par}
    \vspace{1.5cm}
    \rule{0.5\textwidth}{0.4pt}\par
    \vspace{1cm}
    {\large Ken Overton \par}
    \vspace{0.3cm}
    {\normalsize \textit{KJO Liquid Rocket Motor Control System} \par}
    \vspace{1.5cm}
    \begin{tabular}{ll}
      \textbf{Version:} & #3 \\[4pt]
      \textbf{Date:}    & #4 \\
    \end{tabular}
    \vspace{1.5cm}\par
    \begin{minipage}{0.75\textwidth}
      \centering
      \normalsize #5
    \end{minipage}
    \vfill
    {\small\color{KJOgray} \textcopyright\ 2025--2026 Ken Overton.
     All rights reserved. \par}
  \end{titlepage}%
}

\newcommand{\docversion}{0.7}
\newcommand{\docdate}{September 2026}
\lhead{\textsc{KJO GSE Management Unit}}
\rhead{\textsc{GSEMU Reference} --- v\docversion}
\cfoot{\thepage\ of \pageref{LastPage}}
\hypersetup{pdftitle={KJO GSEMU Reference}, pdfauthor={Ken Overton}}

\begin{document}
% KJO_tikz_styles.tex
% Shared TikZ node and edge styles for all KJO block diagrams.
% \input this file inside the document body before any tikzpicture environments.

\tikzset{
  % ── Hardware block (peripheral chip / PCB component) ──────────────────────
  hwblock/.style={
    rectangle, draw=KJOgray, line width=1pt,
    fill=KJOgray!10, rounded corners=4pt,
    minimum width=2.6cm, minimum height=0.8cm,
    font=\small, align=center, inner sep=4pt,
  },
  % Sub-variant: MCU (larger, blue border)
  mcublock/.style={
    hwblock,
    draw=KJOblue, fill=KJOblue!8,
    minimum width=3.0cm, minimum height=1.0cm,
    font=\small\bfseries,
  },
  % Sub-variant: power / battery
  powerblock/.style={
    hwblock,
    draw=KJOorange, fill=KJOorange!8,
  },
  % ── Software module box ────────────────────────────────────────────────────
  swmodule/.style={
    rectangle, draw=KJOblue, line width=0.8pt,
    fill=KJOblue!6, rounded corners=3pt,
    minimum width=2.8cm, minimum height=0.7cm,
    font=\small, align=center, inner sep=3pt,
  },
  % Sub-variant: shared library
  sharedlib/.style={
    swmodule,
    draw=KJOpurple, fill=KJOpurple!6,
    dashed,
  },
  % ── Bus / link annotation ──────────────────────────────────────────────────
  busline/.style={
    draw=KJOgray, line width=1.5pt,
    -{Stealth[length=6pt,width=4pt]},
  },
  bidbusline/.style={
    draw=KJOgray, line width=1.5pt,
    {Stealth[length=6pt,width=4pt]}-{Stealth[length=6pt,width=4pt]},
  },
  radioline/.style={
    draw=KJOblue!60, line width=1.5pt, dashed,
    {Stealth[length=6pt,width=4pt]}-{Stealth[length=6pt,width=4pt]},
  },
  gpioline/.style={
    draw=KJOgreen!60!black, line width=1.2pt,
    -{Stealth[length=5pt,width=4pt]},
  },
  % ── Flowchart shapes ──────────────────────────────────────────────────────
  flowproc/.style={
    rectangle, draw=KJOblue!70, line width=0.8pt,
    fill=KJOblue!8, rounded corners=3pt,
    minimum width=3.2cm, minimum height=0.7cm,
    font=\small, align=center, inner sep=3pt,
  },
  flowdec/.style={
    diamond, draw=KJOorange, line width=0.8pt,
    fill=KJOorange!8,
    minimum width=2.2cm, minimum height=0.9cm,
    font=\small, align=center, inner sep=1pt,
  },
  flowterminal/.style={
    rounded rectangle, draw=KJOblue, line width=1pt,
    fill=KJOblue!12,
    minimum width=2.8cm, minimum height=0.7cm,
    font=\small\bfseries, align=center, inner sep=3pt,
  },
  flowarrow/.style={
    draw=KJOgray, line width=0.9pt,
    -{Stealth[length=5pt,width=4pt]},
  },
  % ── Propellant schematic styles ────────────────────────────────────────────
  % Generic pipe line
  pipe/.style={
    draw=black, line width=1.8pt,
  },
  pipearrow/.style={
    pipe,
    -{Stealth[length=7pt,width=5pt]},
  },
  % Oxidizer (N2O) — blue
  oxpipe/.style={
    draw=KJOblue, line width=1.8pt,
  },
  oxpipearrow/.style={
    oxpipe,
    -{Stealth[length=7pt,width=5pt]},
  },
  % Fuel — red
  fuelpipe/.style={
    draw=KJOred, line width=1.8pt,
  },
  fuelpipearrow/.style={
    fuelpipe,
    -{Stealth[length=7pt,width=5pt]},
  },
  % GSE fill line — green
  gsepipe/.style={
    draw=KJOgreen!70!black, line width=1.8pt,
  },
  gsepipearrow/.style={
    gsepipe,
    -{Stealth[length=7pt,width=5pt]},
  },
  % Pressurization — orange dashed
  presspipe/.style={
    draw=KJOorange, line width=1.4pt, dashed,
  },
  presspipearrow/.style={
    presspipe,
    -{Stealth[length=6pt,width=4pt]},
  },
  % Tank (tall rectangle, rounded corners)
  tank/.style={
    rectangle, draw=black, line width=1.5pt,
    rounded corners=4pt, fill=white,
    minimum width=1.8cm, minimum height=2.4cm,
    font=\small\bfseries, align=center, inner sep=4pt,
  },
  n2otank/.style={
    tank, draw=KJOblue, fill=KJOblue!8,
  },
  fueltank/.style={
    tank, draw=KJOred, fill=KJOred!8,
  },
  srctank/.style={
    tank, draw=KJOgreen!70!black, fill=KJOgreen!6,
  },
  % Valve box (servo-actuated)
  valvebox/.style={
    rectangle, draw=black, line width=1pt,
    fill=white, minimum width=2.2cm, minimum height=0.65cm,
    font=\footnotesize\bfseries, align=center, inner sep=3pt,
  },
  gseval/.style={
    valvebox, draw=KJOgreen!70!black, fill=KJOgreen!8,
  },
  emuval/.style={
    valvebox, draw=KJOred, fill=KJOred!6,
  },
  % Component (generic inline component)
  component/.style={
    rectangle, draw=KJOgray, line width=0.8pt,
    fill=white!96!gray, minimum width=2.2cm, minimum height=0.6cm,
    font=\footnotesize, align=center, inner sep=3pt,
  },
  % Instrument callout circle (PT, TT, TC)
  instrument/.style={
    circle, draw=KJOgray, line width=0.8pt,
    fill=white, minimum size=0.6cm,
    font=\tiny\bfseries, align=center, inner sep=0pt,
  },
  % Sensor lead line
  sensorlead/.style={
    draw=KJOgray, line width=0.7pt,
  },
  % Boundary box (dashed outline for GSE / Rocket regions)
  gseboundary/.style={
    draw=KJOgreen!50, line width=1pt, dashed,
    rounded corners=6pt,
  },
  rocketboundary/.style={
    draw=KJOred!40, line width=1pt,
    rounded corners=6pt,
  },
  % Sequence diagram lifeline
  lifeline/.style={
    draw=KJOgray, line width=0.6pt, dashed,
  },
  % Sequence diagram message arrow
  seqarrow/.style={
    draw=KJOblue, line width=0.9pt,
    -{Stealth[length=5pt,width=4pt]},
  },
  % State node (FSM)
  statenode/.style={
    circle, draw=KJOblue, line width=1pt,
    fill=KJOblue!10,
    minimum size=1.2cm,
    font=\small\bfseries, align=center,
  },
  statetransition/.style={
    draw=KJOgray, line width=0.9pt,
    -{Stealth[length=5pt,width=4pt]},
    font=\footnotesize,
  },
}


\KJOtitlepage%
  {GSE Management Unit (GSEMU)}%
  {Hardware and Software Reference}%
  {\docversion}{\docdate}%
  {This document describes the hardware, software architecture, Fill valve
   control, and GPIO handshake protocol of the GSE Management Unit.
   The GSEMU is mounted on the launch pad and controls the propellant fill
   valve and the umbilical quick-release servo.}

\tableofcontents\newpage

%═══════════════════════════════════════════════════════════════════════════════
\section{Hardware Overview}
%═══════════════════════════════════════════════════════════════════════════════

\subsection{Component Summary}

\begin{center}
\begin{tabularx}{\textwidth}{@{}lp{6.2cm}lY@{}}
\toprule
Component & Part & Interface & Notes \\
\midrule
Microcontroller & Adafruit Feather M0 (ATSAMD21G18A, 48~MHz)   & ---     & \\
OLED Display    & SH1107, 128$\times$64, 1.12\textquotedbl       & I\textsuperscript{2}C & \\
GPIO Expander   & MCP23X17 (16 I/O)                             & I\textsuperscript{2}C & Buttons, AUX\_IO \\
PWM Servo Wing  & Adafruit 8-channel wing (PCA9685)             & I\textsuperscript{2}C & Fill valve servo (ch.~6) and QR release servo (ch.~5) \\
ADC             & ADS1015 (12-bit, 4-ch)                       & I\textsuperscript{2}C & Battery voltage, AUX/LCO sense \\
CAN Controller  & MCP2515 (Adafruit CAN FeatherWing)            & SPI     & Link to EMU/LCMU \\
Real-Time Clock & PCF8523                                       & I\textsuperscript{2}C & \\
SD Card         & Standard SPI SD socket                        & SPI     & \\
Battery         & 2S LiPo (7.4~V nominal)                       & ---     & Read via ADS1015, resistor-divider sense \\
\bottomrule
\end{tabularx}
\end{center}

\subsection{MCP23X17 GPIO Expander Pin Map}

\begin{center}
\begin{tabularx}{\textwidth}{@{}llY@{}}
\toprule
Pin Name & Direction & Function \\
\midrule
Button A (QR release)  & IN  & Hold-to-release QR latch (servo open while held, hold when released) \\
Button B (fill close)  & IN  & Locally command Fill valve closed \\
Button C (fill open)   & IN  & Locally command Fill valve open \\
\sig{AUX\_IO\_1} (\sig{GSEMU\_LINK\_SENSE\_EIO}) & IN (INPUT\_PULLUP) & Umbilical link-disconnect sense; wired to EMU chassis GND (LOW = connected, HIGH = separated). Fill valve CMD/STATUS formerly carried here has moved to CAN. \\
\sig{AUX\_IO\_2}                                  & --- & Retired/unused (formerly \sig{FILL\_VALVE\_STATUS\_EIO}). \\
\sig{AUX\_IO\_3} (\sig{QR\_SENSE\_EIO})           & IN (INPUT\_PULLUP) & Umbilical QR physical-connection sense; wired to EMU chassis GND (LOW = connected, HIGH = separated). QR release itself is commanded over CAN (\cmd{CAN\_QR\_RELEASE}), not this line. \\
\sig{AUX\_IO\_4}                                  & --- & Retired/unused (formerly \sig{RELEASE\_STATE\_EIO}; its sense function consolidated onto \sig{QR\_SENSE\_EIO}). \\
\sig{AUX\_IO\_5}                                  & --- & Retired/unused (formerly \sig{LAUNCH\_ENABLE\_EIO}). \\
\sig{AUX\_IO\_6}                                  & --- & Retired/unused (formerly \sig{REMOTE\_START\_EIO}). \\
\bottomrule
\end{tabularx}
\end{center}

\subsection{Hardware Block Diagram}

\begin{figure}[H]
\centering
\resizebox{\textwidth}{!}{%
\begin{tikzpicture}[node distance=0.9cm and 1.8cm]
  \node[mcublock, minimum width=3.0cm, minimum height=1.1cm] (mcu)
    {Feather M0\\(ATSAMD21G18A)};
  % I2C left
  \node[hwblock, left=3.6cm of mcu, yshift= 2.4cm, minimum width=3.2cm] (oled) {SH1107 OLED};
  \node[hwblock, left=3.6cm of mcu, yshift= 1.2cm, minimum width=3.2cm] (gpio) {MCP23X17\\GPIO Exp.};
  \node[hwblock, left=3.6cm of mcu, yshift= 0.0cm, minimum width=3.2cm] (pwm)  {PCA9685\\Servo Wing};
  \node[hwblock, left=3.6cm of mcu, yshift=-1.2cm, minimum width=3.2cm] (adc)  {ADS1015 ADC};
  \node[hwblock, left=3.6cm of mcu, yshift=-2.4cm, minimum width=3.2cm] (rtc)  {PCF8523 RTC};
  % Buttons -- adjacent to the GPIO expander (its output), clear of the
  % mcu/left-column bus lines and the mcu/battery power line below.
  \node[hwblock, left=1.0cm of gpio, minimum width=2.4cm] (btn) {Panel\\Buttons};
  % SPI right
  \node[hwblock, right=2.3cm of mcu, yshift=0.6cm, minimum width=3.0cm] (sd)   {SD Card};
  \node[hwblock, right=2.3cm of mcu, yshift=-0.6cm, minimum width=3.0cm] (can)  {MCP2515\\CAN FeatherWing};
  % Power -- straight down from mcu; nothing else occupies this column.
  \node[powerblock, below=3.0cm of mcu, minimum width=3.0cm] (bat) {2S LiPo Battery};
  % I2C bus
  \foreach \nd in {oled,gpio,pwm,adc,rtc}
    \draw[bidbusline] (mcu.west |- \nd) -- (\nd.east)
      node[midway,above,font=\tiny]{I\textsuperscript{2}C};
  % SPI
  \foreach \nd in {sd,can}
    \draw[bidbusline] (mcu.east |- \nd) -- (\nd.west)
      node[midway,above,font=\tiny]{SPI};
  % GPIO expander → buttons (simple adjacency, no crossing)
  \draw[busline] (gpio.west) -- (btn.east);
  % Power -- straight vertical, clear of every other node.
  \draw[busline,draw=KJOorange] (bat.north) -- (mcu.south)
    node[midway,right,font=\tiny]{7.4V};
\end{tikzpicture}%
}
\caption{GSEMU hardware block diagram.}
\end{figure}

%═══════════════════════════════════════════════════════════════════════════════
\section{Software Architecture}
%═══════════════════════════════════════════════════════════════════════════════

\subsection{Module Summary}

\begin{center}
\begin{tabularx}{\textwidth}{@{}llY@{}}
\toprule
Module & Location & Responsibility \\
\midrule
\sw{main.cpp}            & local  & Setup, loop, button handling, CAN command dispatch (\sw{Check\_CAN()}), deferred fill-move response state machine, battery/LCO polling, QR test-release timing \\
\sw{KJO\_Valve}          & local  & Base valve: servo drive, pot feedback (Fill valve open/close is CAN-commanded; see §4) \\
\sw{KJO\_QR\_Servo}      & local  & Servo-only actuator base class (time-based position) \\
\sw{KJO\_QR\_Slave}      & local  & QR release actuator: latches \cmd{CAN\_QR\_RELEASE}, reads connection-sense (\sig{AUX\_IO\_3}), local override \\
\sw{KJO\_GPIO}           & local  & MCP23X17 pin assignments, QR servo constants \\
\sw{KJO\_Analog}         & local  & ADS1015 channel assignments, LiPo scale, AUX/LCO threshold \\
\sw{KJO\_GSE\_Display}   & local  & SD log-file finder (\sw{Find\_Available\_File()}) only -- OLED display moved to shared \sw{KJO\_Status\_Display} \\
\sw{KJO\_Logging}        & local  & SD log file constants \\
\sw{KJO\_System\_Config} & shared & Version, platform constants \\
\sw{KJO\_CAN\_Command\_Defs} & shared & CAN command codes, packet structures, node IDs \\
\sw{KJO\_CAN\_Client}    & shared & CAN transport (send/poll/request-with-timeout) \\
\sw{KJO\_Valve\_Timing}  & shared & Fill valve move-timeout constant (shared with EMU) \\
\sw{KJO\_LCO\_Sense}     & shared & AUX/LCO ADS1015 channel + threshold (shared with EMU) \\
\sw{KJO\_Status} / \sw{KJO\_Status\_Display} & shared & Tagged Serial/SD logging (\sw{Log\_Message()}) and scrolling OLED status (\sw{Report\_Status()}) \\
\bottomrule
\end{tabularx}
\end{center}

\subsection{Class Hierarchy}

\sw{Fill\_Valve} is a plain \texttt{Valve} instance -- there is no derived
class for it. The \sw{Valve} base class provides servo drive and
potentiometer feedback; \sw{open()}/\sw{close()} are called directly from
\sw{Check\_CAN()} (see §4).

\texttt{QR\_Servo} (base) $\longleftarrow$ \texttt{QR\_Slave} (derived).

\sw{QR\_Servo} provides servo drive via the PCA9685 with time-based motion
completion detection (no potentiometer).  \sw{QR\_Slave} adds a permanent
release-commanded latch (set by \cmd{CAN\_QR\_RELEASE}), physical
separation detection via \sig{AUX\_IO\_3} (\sig{QR\_SENSE\_EIO}), and a
local Button~A override (\sw{localRelease()} / \sw{localHold()}).

\subsection{Main Loop Control Flow}

\begin{figure}[H]
\centering
\begin{tikzpicture}[node distance=0.5cm]
  \node[flowterminal] (start)  {loop()};
  \node[flowproc, below=of start]  (qr)    {QR\_Release.update()};
  \node[flowproc, below=of qr]     (intl)  {Fill Valve safety interlock\\(inline; closes on QR separation)};
  \node[flowproc, below=of intl]   (btn)   {Check\_Buttons()};
  \node[flowproc, below=of btn]    (bat)   {Check\_Battery()};
  \node[flowproc, below=of bat]    (lco)   {Check\_LCO\_Watch()};
  \node[flowproc, below=of lco]    (fv)    {Check\_Fill\_Valve\_Move()};
  \node[flowproc, below=of fv]     (qrt)   {Check\_QR\_Test\_Release()};
  \node[flowproc, below=of qrt]    (pres)  {Check\_CAN\_Presence()};
  \node[flowproc, below=of pres]   (can)   {Check\_CAN()};
  \draw[flowarrow] (start)--(qr);
  \draw[flowarrow] (qr)--(intl);
  \draw[flowarrow] (intl)--(btn);
  \draw[flowarrow] (btn)--(bat);
  \draw[flowarrow] (bat)--(lco);
  \draw[flowarrow] (lco)--(fv);
  \draw[flowarrow] (fv)--(qrt);
  \draw[flowarrow] (qrt)--(pres);
  \draw[flowarrow] (pres)--(can);
  \draw[flowarrow] (can.south) -- ++(0,-0.4) -- ++(4.0,0)
    -- ++(0,10.0) -- (start.east);
\end{tikzpicture}
\caption{GSEMU main loop control flow.}
\end{figure}

%═══════════════════════════════════════════════════════════════════════════════
\section{CAN Command Table}
%═══════════════════════════════════════════════════════════════════════════════

\sw{Check\_CAN()} (called every \sw{loop()}) is GSEMU's sole CAN command
dispatcher. Every command it handles:

\begingroup
\footnotesize
\begin{longtable}{@{}p{3.4cm}p{2.3cm}p{2.1cm}p{6.4cm}@{}}
\toprule
Command & Sent by & Response & Notes \\
\midrule
\cmd{CAN\_\allowbreak GET\_\allowbreak HEALTH\_\allowbreak STATUS} & Any node & uint16 bitfield & SD/RTC init-failure bits plus link-sense-lost \\
\cmd{CAN\_\allowbreak PING} & Any node & ack & Boot-time presence check \\
\cmd{CAN\_\allowbreak ARM\_\allowbreak LCO\_\allowbreak WATCH} & EMU & ack & Arms/disarms \sw{Check\_\allowbreak LCO\_\allowbreak Watch()}'s AUX-channel polling \\
\cmd{CAN\_\allowbreak QR\_\allowbreak RELEASE} & EMU & none (fire-and-forget) & Permanent latch via \sw{commandRelease()}; production launch only \\
\cmd{CAN\_\allowbreak QR\_\allowbreak TEST\_\allowbreak RELEASE} & RCU\_\allowbreak UNIT\_\allowbreak TEST & ack & Reversible, self-reverting test release (see §5) \\
\cmd{CAN\_\allowbreak OPEN\_\allowbreak FILL\_\allowbreak VALVE} & EMU or RCU\_\allowbreak UNIT\_\allowbreak TEST & deferred ack + position & See §4 \\
\cmd{CAN\_\allowbreak CLOSE\_\allowbreak FILL\_\allowbreak VALVE} & EMU or RCU\_\allowbreak UNIT\_\allowbreak TEST & deferred ack + position & See §4 \\
\cmd{CAN\_\allowbreak GET\_\allowbreak GSEMU\_\allowbreak BATTERY} & EMU & float volts & \\
\cmd{CAN\_\allowbreak GET\_\allowbreak FILL\_\allowbreak VALVE\_\allowbreak POSITION} & RCU\_\allowbreak UNIT\_\allowbreak TEST & float \% & Bench-test only; see §4.3 \\
\cmd{CAN\_\allowbreak GET\_\allowbreak LCO\_\allowbreak STATE} & EMU & float (0/1) & Read-only AUX/LCO sense; does not touch \texttt{lco\_\allowbreak watch\_\allowbreak armed} \\
\bottomrule
\end{longtable}
\endgroup

Any command not listed above (e.g. the LCMU-side TARE/weight/thrust
commands) is not addressed to GSEMU and falls through unhandled by design.

%═══════════════════════════════════════════════════════════════════════════════
\section{Fill Valve Control}
%═══════════════════════════════════════════════════════════════════════════════

\subsection{CAN-Commanded Operation}

The wired \sig{AUX\_IO\_1}/\sig{AUX\_IO\_2} CMD/STATUS handshake this
section formerly described has been retired; there is no \sw{Slave\_Valve}
class in the current firmware. \sw{Fill\_Valve} is a plain \sw{Valve}
instance, opened/closed directly on EMU's CAN command. GSEMU has no fill
state machine of its own -- EMU owns tare/poll/target tracking for the
whole fill session.

\begin{sloppypar}
\begin{enumerate}[nosep]
  \item \sw{Check\_CAN()} receives \cmd{CAN\_OPEN\_FILL\_VALVE} or
        \cmd{CAN\_CLOSE\_FILL\_VALVE} and calls \sw{Begin\_Fill\_Move()},
        which issues \sw{Fill\_Valve.open()}/\sw{close()} immediately.
  \item The CAN response is \emph{deferred} -- \sw{Check\_Fill\_Valve\_Move()}
        (called every \sw{loop()}) sends it once \sw{Fill\_Valve.isStopped()}
        confirms the move, or \texttt{MOVE\_TIMEOUT}+100~ms elapses,
        reporting the resulting position either way.
  \item This same mechanism is also reachable directly via
        RCU\_UNIT\_TEST's LoRa \cmd{OPEN\_FILL\_VALVE}/\cmd{CLOSE\_FILL\_VALVE}
        commands (distinct identifiers from, but relayed by EMU as,
        \cmd{CAN\_OPEN\_FILL\_VALVE}/\cmd{CAN\_CLOSE\_FILL\_VALVE}), for
        exercising the valve without a full fill session.
\end{enumerate}
\end{sloppypar}

\subsection{Local Button Override}

Three panel buttons on the MCP23X17 provide local control during ground
operations. Filed here because Buttons~B/C act on the Fill valve, this
section's subject -- Button~A operates the QR release servo (§5's
subject) and is cross-referenced there too:

\begin{center}
\begin{tabularx}{\textwidth}{@{}lY@{}}
\toprule
Button & Action \\
\midrule
A (QR Release) & \textbf{Hold-to-release}: calls \sw{localRelease()} on press; servo opens
                 while button held.  Calls \sw{localHold()} on release; servo returns to
                 hold position.  Event logged to SD card. \\
B (Fill Close) & Drive Fill valve servo to fully closed position (falling-edge triggered). \\
C (Fill Open)  & Drive Fill valve servo to fully open position (falling-edge triggered). \\
\bottomrule
\end{tabularx}
\end{center}

\begin{warningbox}
Button~A activates the umbilical quick-release servo.  The latch opens for
as long as the button is held, and releases pressurised
N\textsubscript{2}O when the umbilical separates.  Ensure all personnel are
clear and the pad is safe before pressing Button~A.
Buttons~B and~C bypass EMU command authority; use for ground testing only.
\end{warningbox}

\subsection{Fill Valve Position Query}

\begin{sloppypar}
\cmd{CAN\_GET\_FILL\_VALVE\_POSITION} is an RCU\_UNIT\_TEST-only bench-test
command (production traffic never sends it) that polls the Fill valve's
live position at any time, answered immediately with
\sw{Fill\_Valve.getPositionPercent()} -- independent of any in-progress
open/close move.
\end{sloppypar}

%═══════════════════════════════════════════════════════════════════════════════
\section{Umbilical Quick-Release (\sw{QR\_Slave})}
%═══════════════════════════════════════════════════════════════════════════════

\subsection{Overview}

\begin{sloppypar}
The umbilical quick-release servo physically separates the umbilical
connector from the rocket airframe at launch. The GSEMU drives the servo;
release is commanded over CAN (\cmd{CAN\_QR\_RELEASE}, fire-and-forget from
EMU) rather than a dedicated wire. Release is one-shot: once
\sw{commandRelease()} is called, the servo opens and stays open -- there
is no remote command back to ``hold.'' A spring-loaded latch, once
triggered, cannot be electrically un-released; only physical re-latching
(and/or the local front-panel override below) changes it.
\end{sloppypar}

\subsection{Hardware}

The QR servo is driven by a channel on the PCA9685 Servo Wing. Unlike the
Fill valve servo, there is no potentiometer feedback; position is determined
entirely by elapsed time (\sw{QR\_SERVO\_MOVE\_MS}).

One MCP23X17 expansion I/O line is used for connection sensing:

\begin{itemize}[nosep]
  \item \sig{AUX\_IO\_3} (\sig{QR\_SENSE\_EIO}) --- INPUT\_PULLUP, wired to
        EMU chassis GND. LOW = umbilical physically connected (EMU GND
        pulls it LOW); HIGH = umbilical has separated.
\end{itemize}

\begin{warningbox}
\texttt{QR\_SERVO\_PWM\_HOLD} and \texttt{QR\_SERVO\_PWM\_OPEN} were
bench-calibrated in March 2026 (superseding earlier placeholder values
seeded from the Fill valve's calibration). Re-verify servo travel
direction and endpoint positions on the bench after any hardware change
before connecting the umbilical.
\end{warningbox}

\subsection{Servo Behaviour}

\sw{QR\_Slave::update()} polls every \sw{loop()} iteration:

\begin{itemize}[nosep]
  \item Local override active OR release commanded, and servo not already
        opening/open $\Rightarrow$ \sw{openServo()}.
  \item Otherwise, and servo not already holding/held $\Rightarrow$
        \sw{holdServo()}.
\end{itemize}

Physical separation (state line LOW $\to$ HIGH) is detected and latched;
\sw{isSeparated()} returns true once separation has been confirmed.
\sw{isCurrentlyConnected()} instead returns the live, non-latching state --
used where behaviour must resume normally on reconnect (e.g. the Fill
Valve safety interlock in \sw{main.cpp}), unlike the permanently-latched
\sw{isSeparated()}.

\subsection{Local Button A Override}

Button~A on the GSEMU front panel engages a local override that bypasses
the release-commanded latch and forces the servo open, independent of
whether CAN release has been commanded. \sw{Check\_Buttons()} calls:

\begin{itemize}[nosep]
  \item \sw{QR\_Release.localRelease()} on button press (falling edge).
  \item \sw{QR\_Release.localHold()} on button release (rising edge).
\end{itemize}

The override is cleared on button release; the servo then follows the
release-commanded latch (open if \sw{CAN\_QR\_RELEASE} was ever received,
hold otherwise) on the next \sw{loop()} iteration.

\subsection{QR\_Slave API}

\begin{center}
\begin{tabularx}{\textwidth}{@{}lY@{}}
\toprule
Method & Action \\
\midrule
\sw{begin()}                & Configure the state pin as INPUT\_PULLUP; command servo to HOLD \\
\sw{update()}                & Drive servo per local override / release-commanded latch; detect separation \\
\sw{commandRelease()}        & Latches the release command (permanent -- never cleared); called from \sw{Check\_CAN()}'s \cmd{CAN\_QR\_RELEASE} case \\
\sw{isSeparated()}           & Latched: true once the state line has gone HIGH (physical separation confirmed) \\
\sw{isCurrentlyConnected()}  & Live, non-latching: true while the state line currently reads LOW \\
\sw{localRelease()}          & Set local override; servo opens on next \sw{update()} \\
\sw{localHold()}             & Clear local override; servo follows the release-commanded latch on next \sw{update()} \\
\bottomrule
\end{tabularx}
\end{center}

%═══════════════════════════════════════════════════════════════════════════════
\section{OLED Display}
%═══════════════════════════════════════════════════════════════════════════════

The SH1107 OLED provides local status when no serial console is available.
The scrolling status display itself lives in the shared
\sw{KJO\_Status\_Display} module (used identically by EMU, GSEMU, and
LCMU), not in \sw{KJO\_GSE\_Display}:

\begin{itemize}[nosep]
  \item \sw{Report\_Status(display, tag, message, delay\_flag)}: logs the
        message (Serial + SD, timestamped, via \sw{Log\_Message()}), then
        redraws the scrolling display showing \texttt{"[TAG] message"}
        (no timestamp on-screen) in a fixed \texttt{MAX\_LINES}-line
        circular buffer, then delays 1~s if \texttt{delay\_flag} is true.
\end{itemize}

\sw{KJO\_GSE\_Display}'s only remaining responsibility is
\sw{Find\_Available\_File()}, which scans the SD card root for the next
unused \texttt{Log\_N.txt} index at boot.

%═══════════════════════════════════════════════════════════════════════════════
\section{SD Logging}
%═══════════════════════════════════════════════════════════════════════════════

Log files are stored in the root of the SD card as \texttt{Log\_N.txt}
(N is auto-incremented to avoid overwriting existing files).
Each record is time-stamped and tagged (via the shared \sw{Log\_Message()}),
recording button presses, CAN command dispatch, Fill valve and QR release
events, and other system events -- a generic tagged status/event log, not
a fixed set of transition types.

%═══════════════════════════════════════════════════════════════════════════════
\section{Parameters \& Flags}
%═══════════════════════════════════════════════════════════════════════════════

This section collects the tunable operational parameters and feature flags
compiled into this firmware image, together with the source file where each
is defined. It is meant as a single reference point when tracing or changing
a timing, threshold, or calibration value. Pin/channel wiring assignments,
CAN command opcodes, and packet byte layouts are documented separately (the
Hardware Overview pin map and the CAN Command Table, §3, above) and are not
repeated here. Parameters marked \emph{(shared)} are defined in
\sw{KJO\_\allowbreak Shared\_\allowbreak Libraries} and must agree with the same constant used on
the other board(s) that consume it.

\subsubsection*{Fill Valve Servo Drive \& Motion Detection}

\begingroup
\footnotesize
\begin{longtable}{@{}p{3.2cm}p{2.4cm}p{6.2cm}p{2.6cm}@{}}
\toprule
Parameter & Value & Description & Source File \\
\midrule
\texttt{SERVO\_\allowbreak FREQUENCY} & 1100~Hz & PWM frequency driven above the standard 50~Hz to compress the 5-turn servo's usable travel into a smaller PWM-count range, improving position resolution. & \texttt{KJO\_\allowbreak Valve.h} \\
\texttt{SERVO\_\allowbreak RANGE\_\allowbreak MIN} & 500~$\mu$s & Servo pulse width at the lower mechanical travel limit. & \texttt{KJO\_\allowbreak Valve.h} \\
\texttt{SERVO\_\allowbreak RANGE\_\allowbreak MAX} & 2500~$\mu$s & Servo pulse width at the upper mechanical travel limit. & \texttt{KJO\_\allowbreak Valve.h} \\
\texttt{SERVO\_\allowbreak FULL\_\allowbreak RANGE\_\allowbreak DEGREES} & 1800\textdegree & Total mechanical rotation of the 5-turn servo across its full pulse-width range. & \texttt{KJO\_\allowbreak Valve.h} \\
\texttt{SERVO\_\allowbreak PINION\_\allowbreak TEETH} & 16 & Tooth count on the servo output pinion, used to compute the servo-to-valve gear ratio. & \texttt{KJO\_\allowbreak Valve.h} \\
\texttt{VALVE\_\allowbreak GEAR\_\allowbreak TEETH} & 60 & Tooth count on the valve actuator gear, used with \texttt{SERVO\_\allowbreak PINION\_\allowbreak TEETH} to compute the gear ratio. & \texttt{KJO\_\allowbreak Valve.h} \\
\texttt{VALVE\_\allowbreak TURN\_\allowbreak DEGREES} & 90\textdegree & Ball-valve rotation from fully open to fully closed. & \texttt{KJO\_\allowbreak Valve.h} \\
\texttt{VALVE\_\allowbreak MOVE\_\allowbreak TIMEOUT\_\allowbreak MS} (alias \texttt{MOVE\_\allowbreak TIMEOUT}) & 2500~ms & Maximum time allowed for a commanded Fill Valve move to complete before it's treated as failed/timed out; shared with EMU since EMU's own CAN-relay timeout for this valve is sized off this value. & \texttt{KJO\_\allowbreak Valve\_\allowbreak Timing.h} (shared) \\
\texttt{VALVE\_\allowbreak POSITION\_\allowbreak AVERAGING\_\allowbreak COUNT} & 4 & Number of ADC reads averaged per position sample to reduce potentiometer noise. & \texttt{KJO\_\allowbreak Valve.h} \\
\texttt{VALVE\_\allowbreak MOVING\_\allowbreak THRESHOLD} & 2~counts & Minimum position-sensor delta between samples required to consider the valve still moving. & \texttt{KJO\_\allowbreak Valve.h} \\
\texttt{VALVE\_\allowbreak POSITION\_\allowbreak DEAD\_\allowbreak BAND} & 15~counts & Position tolerance used to declare the valve ``open,'' ``closed,'' or ``at target.'' & \texttt{KJO\_\allowbreak Valve.h} \\
\texttt{MOVE\_\allowbreak DETECT\_\allowbreak WAIT} & 30~ms & Minimum interval between position samples inside \sw{isMoving()}'s velocity check. & \texttt{KJO\_\allowbreak Valve.h} \\
\texttt{MOVE\_\allowbreak AFTER\_\allowbreak START\_\allowbreak WAIT} & 60~ms & Blocking delay after issuing a new PWM target, giving the servo time to start moving before the first \sw{isMoving()} sample. & \texttt{KJO\_\allowbreak Valve.h} \\
\texttt{USE\_\allowbreak MAPPING} & \texttt{true} & Feature flag selecting angle-geometry-based valve position mapping over simple linear PWM interpolation. & \texttt{KJO\_\allowbreak System\_\allowbreak Config.h} (shared) \\
\end{longtable}
\endgroup

\subsubsection*{Fill Valve Calibration (Travel Endpoints \& Ball-Valve Geometry)}

\textit{Note: the servo/pot channel these values were calibrated against
(PCA9685 channel 6 / ADC pin A1) is itself a temporary diagnostic
reassignment, not the original channel~4/A4 design -- see
\texttt{KJO\_Valve.h} for the bench-fault-isolation rationale. Re-verify
these endpoints if the channel assignment reverts.}

\begingroup
\footnotesize
\begin{longtable}{@{}p{3.2cm}p{2.4cm}p{6.2cm}p{2.6cm}@{}}
\toprule
Parameter & Value & Description & Source File \\
\midrule
\texttt{FILL\_\allowbreak VALVE\_\allowbreak PWM\_\allowbreak OPEN} & 2329 & Servo PWM count corresponding to the Fill Valve's fully-open position. & \texttt{KJO\_\allowbreak Valve.h} \\
\texttt{FILL\_\allowbreak VALVE\_\allowbreak PWM\_\allowbreak CLOSE} & 3918 & Servo PWM count corresponding to the Fill Valve's fully-closed position. & \texttt{KJO\_\allowbreak Valve.h} \\
\texttt{FILL\_\allowbreak VALVE\_\allowbreak POS\_\allowbreak OPEN} & 1310~counts & Potentiometer reading corresponding to the Fill Valve's fully-open position. & \texttt{KJO\_\allowbreak Valve.h} \\
\texttt{FILL\_\allowbreak VALVE\_\allowbreak POS\_\allowbreak CLOSED} & 3004~counts & Potentiometer reading corresponding to the Fill Valve's fully-closed position. & \texttt{KJO\_\allowbreak Valve.h} \\
\texttt{FILL\_\allowbreak VALVE\_\allowbreak BALL\_\allowbreak DIAMETER} & 0.520~in & Ball outer diameter used to compute the angle-based flow-restriction mapping. & \texttt{KJO\_\allowbreak Valve.h} \\
\texttt{FILL\_\allowbreak VALVE\_\allowbreak BORE\_\allowbreak DIAMETER} & 0.275~in & Ball bore diameter used in the same geometric flow-mapping calculation. & \texttt{KJO\_\allowbreak Valve.h} \\
\texttt{FILL\_\allowbreak VALVE\_\allowbreak THROAT\_\allowbreak DIAMETER} & 0.275~in & Valve body throat inner diameter used in the same geometric flow-mapping calculation. & \texttt{KJO\_\allowbreak Valve.h} \\
\end{longtable}
\endgroup

\subsubsection*{Umbilical Quick-Release (QR) Servo}

\begingroup
\footnotesize
\begin{longtable}{@{}p{3.2cm}p{2.4cm}p{6.2cm}p{2.6cm}@{}}
\toprule
Parameter & Value & Description & Source File \\
\midrule
\texttt{QR\_\allowbreak SERVO\_\allowbreak PWM\_\allowbreak HOLD} & 2400 & PWM count for the latch-engaged (held/connected) servo position, bench-calibrated March 2026. & \texttt{KJO\_\allowbreak GPIO.h} \\
\texttt{QR\_\allowbreak SERVO\_\allowbreak PWM\_\allowbreak OPEN} & 2800 & PWM count for the latch-released (open) servo position. & \texttt{KJO\_\allowbreak GPIO.h} \\
\texttt{QR\_\allowbreak SERVO\_\allowbreak MOVE\_\allowbreak MS} & 1000~ms & Time allotted for the QR servo to complete a full stroke between hold and open. & \texttt{KJO\_\allowbreak GPIO.h} \\
\texttt{QR\_\allowbreak TEST\_\allowbreak RELEASE\_\allowbreak DURATION\_\allowbreak MS} & 3000~ms & Bench-test-only dwell time the QR servo stays released before automatically returning to hold (\cmd{QR\_\allowbreak TEST\_\allowbreak RELEASE}, RCU\_\allowbreak UNIT\_\allowbreak TEST); never applied to the permanent production release latch. & \texttt{main.cpp} \\
\end{longtable}
\endgroup

\subsubsection*{LCO Ignition Sense \& Battery Monitoring (Analog Subsystem)}

\begingroup
\footnotesize
\begin{longtable}{@{}p{3.2cm}p{2.4cm}p{6.2cm}p{2.6cm}@{}}
\toprule
Parameter & Value & Description & Source File \\
\midrule
\texttt{AD\_\allowbreak BASE\_\allowbreak SCALE} & 0.00199144777 V/count & ADS1015 volts-per-count scale factor at \texttt{GAIN\_\allowbreak ONE} (4.096~V / 2048~counts); identical for all boards using this gain setting. & \texttt{KJO\_\allowbreak Analog.h} \\
\texttt{GSEMU\_\allowbreak LIPO\_\allowbreak SCALE} & 2.79111 & Voltage-divider scale factor converting raw ADC counts to 2S LiPo battery voltage; a placeholder seeded from EMU's divider, pending measurement of GSEMU's actual divider resistors. & \texttt{KJO\_\allowbreak Analog.h} \\
\texttt{LCO\_\allowbreak SENSE\_\allowbreak THRESHOLD} (alias \texttt{AD\_\allowbreak AUX\_\allowbreak THRESHOLD}) & 251~counts & Digital HIGH/LOW threshold ($\sim$0.5~V post-divider) for the AUX/LCO analog input; bare unsigned comparison, not \texttt{abs()}, since ADS1015 input-protection diodes make a reversed-polarity reading unreliable to detect this way. Shared with EMU's Airstart sense so both boards interpret the same signal identically. & \texttt{KJO\_\allowbreak LCO\_\allowbreak Sense.h} (shared) \\
\texttt{LCO\_\allowbreak SENSE\_\allowbreak ADS\_\allowbreak CHANNEL} (alias \texttt{AD\_\allowbreak AUX\_\allowbreak CHANNEL}) & 0 & ADS1015 channel carrying the AUX/LCO sense signal; primarily a wiring/channel assignment, centralized here only to guarantee EMU and GSEMU agree on which channel carries this cross-board safety signal. & \texttt{KJO\_\allowbreak LCO\_\allowbreak Sense.h} (shared) \\
\texttt{BATTERY\_\allowbreak DISPLAY\_\allowbreak INTERVAL\_\allowbreak MS} & 10000~ms & Interval at which the GSEMU 2S LiPo battery voltage is re-read, displayed, and logged. & \texttt{KJO\_\allowbreak Analog.h} \\
\end{longtable}
\endgroup

\subsubsection*{CAN Communication Timing}

\begingroup
\footnotesize
\begin{longtable}{@{}p{3.2cm}p{2.4cm}p{6.2cm}p{2.6cm}@{}}
\toprule
Parameter & Value & Description & Source File \\
\midrule
\texttt{CAN\_\allowbreak BAUDRATE} & 1{,}000{,}000~bps & CAN bus bit rate, raised from a prior 250~kbps; validated against ISO 11898-2's 1~Mbps/40~m limit for this installation's 16~ft maximum run. & \texttt{KJO\_\allowbreak CAN\_\allowbreak Command\_\allowbreak Defs.h} (shared) \\
\end{longtable}
\endgroup

\subsubsection*{SD Logging}

\begingroup
\footnotesize
\begin{longtable}{@{}p{3.2cm}p{2.4cm}p{6.2cm}p{2.6cm}@{}}
\toprule
Parameter & Value & Description & Source File \\
\midrule
\texttt{LOG\_\allowbreak FILE\_\allowbreak NAME\_\allowbreak BASE} & \texttt{"Log\_\allowbreak "} & Base filename prefix for sequentially-numbered SD log files (e.g.\ \texttt{Log\_\allowbreak 0.txt}), written flat to the SD card root. & \texttt{KJO\_\allowbreak Logging.h} \\
\end{longtable}
\endgroup

\subsubsection*{OLED Boot/Status Scroll Display}

\begingroup
\footnotesize
\begin{longtable}{@{}p{3.2cm}p{2.4cm}p{6.2cm}p{2.6cm}@{}}
\toprule
Parameter & Value & Description & Source File \\
\midrule
\texttt{MAX\_\allowbreak LINES} & 8 & Number of most-recent status lines retained in the scrolling OLED display's circular message buffer. & \texttt{KJO\_\allowbreak Status\_\allowbreak Display.cpp} (shared) \\
\end{longtable}
\endgroup

%═══════════════════════════════════════════════════════════════════════════════
\section{Version History}
%═══════════════════════════════════════════════════════════════════════════════
\begin{center}
\begin{tabularx}{\textwidth}{@{}llY@{}}
\toprule
Version & Date & Notes \\
\midrule
0.1 & March 2026 & Initial release. \\
0.2 & March 2026 & Add umbilical QR release subsystem (QR\_Servo, QR\_Slave, AUX\_IO\_3 protocol). \\
0.3 & March 2026 & Level-based QR protocol (CMD LOW=release, HIGH=hold); add RELEASE\_STATE\_EIO (AUX\_IO\_4),
                   LAUNCH\_ENABLE\_EIO (AUX\_IO\_5), REMOTE\_START\_EIO (AUX\_IO\_6);
                   Button~A reassigned to hold-to-release QR latch;
                   Button~C now opens Fill valve; \sw{QR\_Slave} API adds \sw{localRelease()},
                   \sw{localHold()}, \sw{isSeparated()};
                   \sw{Check\_AUX\_Input()} drives gated Remote Start output. \\
0.4 & March 2026 & SD log files moved to SD card root (\texttt{Log\_N.txt}); subdirectory path
                   constant (\texttt{LOG\_FILE\_FOLDER}) removed from \sw{KJO\_Logging.h} and
                   \sw{KJO\_GSE\_Display}. \\
0.5 & September 2026 & Add new §7 Parameters \& Flags reference section: tunable timing,
                   threshold, and calibration constants across the Fill valve, QR servo,
                   LCO sense, and CAN subsystems, with source-file attribution for each. \\
0.6 & September 2026 & Corrected stale wired-GPIO descriptions left over from the CAN
                   redesign: \S3 Fill Valve Control rewritten (no \sw{Slave\_Valve} class
                   exists; \sw{Fill\_Valve} is CAN-commanded directly from
                   \sw{Check\_CAN()}); \S4 Umbilical Quick-Release (\sw{QR\_Slave})
                   rewritten to describe the permanent \cmd{CAN\_QR\_RELEASE} latch instead
                   of a level-based CMD line; corrected 5 of 6 \sig{AUX\_IO} pin-map rows to
                   match (only the umbilical/link-sense functions survive); fixed the Module
                   Summary and Class Hierarchy (\S2) and Main Loop Control Flow diagram
                   (\S2.3) to match. \\
0.7 & September 2026 & Doc-vs-code audit fixes: resolved a self-contradiction on QR servo
                   PWM calibration (bench-calibrated March 2026, not a placeholder -- the old
                   warning box was itself stale); completed the Main Loop diagram (added the
                   Fill Valve safety interlock and \sw{Check\_CAN()}, both previously missing);
                   added a new \S3 CAN Command Table (all 10 commands \sw{Check\_CAN()}
                   handles, several previously undocumented anywhere); rewrote \S6 OLED
                   Display around the real shared \sw{Report\_Status()} API (the previously
                   documented \sw{scrollMessage()} etc.\ do not exist); corrected the battery
                   chemistry (2S, not 1S) in the Component Summary; added the ADS1015 ADC and
                   MCP2515 CAN controller to the Component Summary and Hardware Block Diagram
                   (both omitted despite being central to the rest of the document); expanded
                   the Module Summary to include all locally- and shared-header modules
                   actually in use; fixed a dangling ``Command Set sections above''
                   cross-reference; clarified that \cmd{CAN\_GET\_FILL\_VALVE\_POSITION} is an
                   RCU\_UNIT\_TEST-only bench command, not normal EMU traffic. \\
\bottomrule
\end{tabularx}
\end{center}

\end{document}
